% !TEX encoding = UTF-8 Unicode
% !TEX program = xelatex
% Updated September 2026; bilingual content synchronized.
\documentclass[margin]{res}
\textwidth=5.2in
\usepackage{fancyhdr}
\pagestyle{plain}
\usepackage{url}
\usepackage{graphicx}
\usepackage{amsmath,amssymb}
\usepackage{fontspec}
\usepackage{xcolor}
\usepackage[hidelinks]{hyperref}
\setmainfont{Hoefler Text}
\setsansfont[Scale=MatchLowercase]{Gill Sans}
\setmonofont[Scale=MatchLowercase]{Andale Mono}
\renewcommand{\labelitemi}{$\bullet$}
\newcommand{\myname}{\textbf{\textcolor{blue}{Mohan Xu}}}

\begin{document}
\enlargethispage{\baselineskip}
\moveleft.5\hoffset\centerline{\large\bf Mohan Xu}
\moveleft.5\hoffset\centerline{DETAILED CURRICULUM VITAE}
\moveleft.5\hoffset\centerline{September 2026}
\moveleft\hoffset\vbox{\hrule width\resumewidth height 1pt}\smallskip

\address{Guanghua School of Management\\
Peking University, Beijing, China}
\address{Email: xu.mohan@stu.pku.edu.cn\\
Website: \href{https://mohan-xu.github.io/}{mohan-xu.github.io}}

\begin{resume}
\section{Education}
\begin{itemize}
  \item Ph.D. in Economics, Peking University\\
  Guanghua School of Management \hfill 2020-2027 (expected)\\
  Doctoral advisor: Yao Tang
  \item Bachelor of Economics, China Foreign Affairs University\\
  School of International Economics \hfill 2016-2020\\
  Undergraduate advisor: Shaojun Fu
  \item Bachelor of Laws, China Foreign Affairs University\\
  Department of International Law \hfill 2016-2020\\
  Undergraduate advisor: Chunyan Zhang
\end{itemize}

\section{Fields of Interest}
Macroeconomics, Monetary Economics, Large Language Models, the Chinese Economy, International Economics

\section{Research Papers}
\textbf{\textit{Working Papers}}
\begin{enumerate}
  \item From Talk to Walk: Fiscal Communication and Asset Prices in China (2026)\\
  \textbf{(Job Market Paper)}\\
  Authors: \myname, Yixin Zhang, Runheng Li and Yao Tang\\
  Summary: Builds an independently annotated knowledge base and combines it with retrieval-augmented generation (RAG) for large language models to identify fiscal policy stances in official Chinese communications and analyze the effects of fiscal expansion on asset prices and the macroeconomy.

  \item Economic Disagreement Across U.S. and Chinese Large Language Models: Evidence from Equity Markets\\
  Authors: \myname, Yuankuan He, Yao Tang and Feng Li\\
  Summary: Collects real-time economic forecasts from U.S. and Chinese large language models, documents their disagreement and a home-market advantage in equity return forecasts, and uses this disagreement to construct forecast combinations.

  \item From Policy Text to Policy Response Functions: Inference with Generated Semantic Partitions\\
  Authors: Pengcheng Ding, Mengxi Yu, \myname, Feng Li and Yao Tang\\
  Summary: Develops generated semantic partitions for structured analysis of policy texts and uses them to infer policy response functions.

  \item The Phillips Curve in the Mind: Experimental Evidence from Chinese Business Executives\\
  Authors: Boning Liu, Yao Tang and \myname\\
  Summary: Uses a survey of Chinese entrepreneurs with an embedded information intervention experiment to show how historical benchmark information revises inflation perceptions, growth expectations, and beliefs about the relationship between the two.

  \item What Explains the Greening of China's Energy Outward Direct Investment? A Host-Country Policy Perspective\\
  Authors: \myname\ and Yao Tang\\
  Summary: Documents the sustained rise in the share of renewable energy projects in China's outward direct investment in energy and analyzes the effects of host-country environmental regulations and financial conditions.
\end{enumerate}

\textbf{\textit{Conference Paper}}
\begin{enumerate}
  \item FMPAF: How Do Fed Chairs Affect the Financial Market?\\
  A Fine-grained Monetary Policy Analysis Framework on Their Language (with Yayue Deng and Yao Tang), \textit{Association for the Advancement of Artificial Intelligence 2024 Workshop on AI in Finance for Social Impact}.
  \begin{itemize}
    \item Based on this research, the team was granted a Chinese invention patent, ``A Method and System for Multimodal Fine-Grained Stance Analysis Based on Large Models'' (Patent No. ZL 2024 1 0159489.0; title translated from Chinese).
  \end{itemize}
\end{enumerate}

\textbf{\textit{Work in Progress}}
\begin{enumerate}
  \item Decoding Fed Speak: Does Transparency in Policy Communication Matter?\\
  (with Pengcheng Ding, Mengxi Yu, Feng Li and Yao Tang)\\
  Project summary: Uses large language models to quantify the clarity of Federal Reserve chairs' communications and assess its effects on financial markets.

  \item Monetary Policy Communication and the Yield Curve\\
  (with Payne Hennigan and Yao Tang)\\
  Project summary: Uses large language models to conduct a structured analysis of Federal Reserve chairs' communication texts, identify their underlying policy rules, and assess their effects on the yield curve.
\end{enumerate}

\section{Conferences and\\Seminars}
Paper Presentations
\begin{itemize} \itemsep -2pt
  \item Asia Meeting of the Econometric Society-China (AMES2026-China)\\
  Hong Kong, China \hfill June 2026\\
  \emph{From Talk to Walk to Stock: How Does Fiscal Communication Move Financial Markets in China?} (with Yixin Zhang, Runheng Li and Yao Tang)

  \item The 13th World Congress of the Econometric Society\\
  Seoul, South Korea \hfill August 2025\\
  \emph{Decoding Fed Speak: Does Transparency in Policy Communication Matter?} (with Pengcheng Ding, Mengxi Yu, Yao Tang and Feng Li)

  \item The 45th International Symposium on Forecasting\\
  Beijing, China \hfill July 2025\\
  Title at presentation: \emph{When Algorithms Speak with Accents: Economic Forecasting Disparities Across Language-Specific LLMs} (with Yuankuan He, Feng Li and Yao Tang)

  \item 2025 Peking-Tsinghua-Waseda Forum on Macroeconomics\\
  Beijing, China \hfill May 2025\\
  \emph{FMPAF: How Do Fed Chairs Affect the Financial Market?} (with Yayue Deng and Yao Tang)

  \item University of International Business and Economics\\
  Beijing, China \hfill January 2025\\
  \emph{Use LLM to Decode Macro-Policy Communications} (with Yao Tang)

  \item The 38th Annual AAAI Conference on Artificial Intelligence\\
  Vancouver, Canada \hfill February 2024\\
  \emph{FMPAF: How Do Fed Chairs Affect the Financial Market?} (with Yayue Deng and Yao Tang)
\end{itemize}

\newpage
\section{Teaching}
Teaching Assistant, Guanghua School of Management, Peking University\\
(Responsibilities included grading assignments and exams, holding office hours, and leading recitation sessions.)
\begin{itemize} \itemsep -2pt
  \item International Economics (undergraduate) \hfill Spring 2025
  \item Financial Econometrics (graduate) \hfill Fall 2024
  \item Selected Frontier Topics in Macroeconomics (graduate) \hfill Fall 2024
  \item Microeconomics (undergraduate) \hfill Fall 2024
  \item Chinese Economy (undergraduate) \hfill Fall 2022
  \item Game Theory (graduate) \hfill Spring 2022
  \item China in the Global Economy (MBA)\\
  \mbox{}\hfill Fall 2021, Fall 2023, Fall 2024
\end{itemize}

\section{Research Projects\\and Policy Studies}
\textbf{Research Assistant}
\begin{itemize} \itemsep -2pt
  \item Professor Yao Tang \hfill 2021-2025\\
  Guanghua School of Management, Peking University
  \item Professor Qinghua Zhang \hfill 2021\\
  Guanghua School of Management, Peking University
\end{itemize}


% Terminology checked against official sources, September 2026:
% Chinese modernization: https://english.www.gov.cn/news/202407/01/content_WS6682ade6c6d0868f4e8e8bf2.html
% NDRC: https://en.ndrc.gov.cn/
% Energy Foundation China: https://www.efchina.org/About-Us-en/Introduction-en
% 15th Five-Year Plan period: https://en.ndrc.gov.cn/policies/202110/t20211027_1301020.html
% The full English titles of the first three projects below are translations,
% not independently verified official project titles. Chinese source titles retained.
\textbf{Research Projects and Policy Research (Participant)}
\begin{itemize} \itemsep 5pt
  \item A Study of Development Goals and an Evaluation Indicator System\\
  for Chinese Modernization\\
  NDRC (Major Research Project) \hfill Oct. 2022-Jan. 2023
%   Developed indicators of international circulation, examined the evolution and determinants of China's opening-up, and proposed benchmark targets for 2035 and 2050 and policy recommendations.
  \item A Study of Economic and Social Development Goals and Indicators\\
  for the 15th Five-Year Plan Period\\
  NDRC \hfill June-Dec. 2024
%   Constructed a set of international trade indicators, estimated sector-level dependence on foreign trade, and conducted cross-country comparisons to identify differences in levels and structure.
  \item Analysis of Optimal Pathways to Carbon Neutrality in the Transport Sector\\
  Energy Foundation China \hfill Apr.-Oct. 2024
%   Applied the Avoid-Shift-Improve framework to model changes in the composition of transport emissions and identify pathways for restructuring the transport system to meet the 2060 emissions reduction target.
\end{itemize}

\section{Non-academic\\Internships}
\begin{itemize} \itemsep -2pt
  \item Enterprise Research Institute \hfill 2021-2022\\
  Development Research Center of the State Council\\
  Policy Research Assistant
  \item Yinhua Fund Management \hfill 2026\\
  Macro Research Team, Fixed Income Research Department\\
  Macro Quantitative Research Intern
\end{itemize}

\section{Services}
\begin{itemize} \itemsep -2pt
  \item Organizer, Macroeconomics and International Economics Reading Group\\
  Guanghua School of Management, Peking University \hfill 2024
  \item Conference Volunteer, The 2025 International Symposium on Forecasting (ISF) Annual Meeting \hfill 2025
  \item Discussant, The 2025 Peking-Tsinghua-Waseda Forum on Macroeconomics \hfill 2025
\end{itemize}

\section{Scholarships\\and Honors}
\begin{itemize} \itemsep -2pt
  \item Peking University, Doctoral Assistantship Scholarship \hfill 2022-2026
  \item Peking University, Graduate Academic Scholarship \hfill 2020-2022
\end{itemize}


\section{Programming\\and Languages}
\begin{itemize} \itemsep -2pt
  \item Programming: Python, STATA, Matlab, SQL
  \item Languages: Chinese (native); English (fluent; TEM8: Excellent)
\end{itemize}

\end{resume}
\end{document}
